\hypertarget{conditionnelles-et-ruxe9cursivituxe9}{%
\section{Conditionnelles et
récursivité}\label{conditionnelles-et-ruxe9cursivituxe9}}

\hypertarget{division-entiuxe8re-et-modulo}{%
\subsection{Division entière et
modulo}\label{division-entiuxe8re-et-modulo}}

L\textquotesingle opérateur de division entière (\emph{floor division}),
\texttt{//}, divise deux nombres et arrondit le résultat à
l\textquotesingle entier inférieur. Par exemple, supposons que le temps
d\textquotesingle exécution d\textquotesingle un programme soit de 105
minutes. Vous pourriez vouloir savoir combien d\textquotesingle heures
et de minutes cela représente. La division entière renvoie le nombre
entier d\textquotesingle heures, en ignorant le reste :

\begin{Shaded}
\begin{Highlighting}[]
\OperatorTok{\textgreater{}\textgreater{}\textgreater{}}\NormalTok{ minutes }\OperatorTok{=} \DecValTok{105}
\OperatorTok{\textgreater{}\textgreater{}\textgreater{}}\NormalTok{ minutes }\OperatorTok{//} \DecValTok{60}
\DecValTok{1}
\end{Highlighting}
\end{Shaded}

Pour obtenir le reste, vous pouvez soustraire une heure, ou utiliser
l\textquotesingle opérateur modulo, \texttt{\%}, qui divise deux nombres
et renvoie le reste :

\begin{Shaded}
\begin{Highlighting}[]
\OperatorTok{\textgreater{}\textgreater{}\textgreater{}}\NormalTok{ reste }\OperatorTok{=}\NormalTok{ minutes }\OperatorTok{\%} \DecValTok{60}
\OperatorTok{\textgreater{}\textgreater{}\textgreater{}}\NormalTok{ reste}
\DecValTok{45}
\end{Highlighting}
\end{Shaded}

L\textquotesingle opérateur modulo est très utile. Par exemple, vous
pouvez vérifier si un nombre est divisible par un autre : si
\texttt{x\ \%\ y} donne zéro, alors \texttt{x} est divisible par
\texttt{y}.

Vous pouvez également extraire le chiffre le plus à droite
d\textquotesingle un nombre. Par exemple, \texttt{x\ \%\ 10} renvoie le
chiffre des unités de \texttt{x} (en base 10).

\hypertarget{expressions-booluxe9ennes}{%
\subsection{Expressions booléennes}\label{expressions-booluxe9ennes}}

Une expression booléenne est une expression qui est soit vraie
(\texttt{True}), soit fausse (\texttt{False}). Les valeurs suivantes
sont de type \texttt{bool} :

\begin{Shaded}
\begin{Highlighting}[]
\OperatorTok{\textgreater{}\textgreater{}\textgreater{}} \BuiltInTok{type}\NormalTok{(}\VariableTok{True}\NormalTok{)}
\OperatorTok{\textless{}}\KeywordTok{class} \StringTok{\textquotesingle{}bool\textquotesingle{}}\OperatorTok{\textgreater{}}
\OperatorTok{\textgreater{}\textgreater{}\textgreater{}} \BuiltInTok{type}\NormalTok{(}\VariableTok{False}\NormalTok{)}
\OperatorTok{\textless{}}\KeywordTok{class} \StringTok{\textquotesingle{}bool\textquotesingle{}}\OperatorTok{\textgreater{}}
\end{Highlighting}
\end{Shaded}

L\textquotesingle opérateur \texttt{==} compare deux opérandes et
produit \texttt{True} s\textquotesingle ils sont égaux, et
\texttt{False} sinon. Les autres opérateurs de comparaison relationnelle
sont :

\begin{Shaded}
\begin{Highlighting}[]
\NormalTok{x }\OperatorTok{!=}\NormalTok{ y               }\CommentTok{\# x n\textquotesingle{}est pas égal à y}
\NormalTok{x }\OperatorTok{\textgreater{}}\NormalTok{ y                }\CommentTok{\# x est strictement supérieur à y}
\NormalTok{x }\OperatorTok{\textless{}}\NormalTok{ y                }\CommentTok{\# x est strictement inférieur à y}
\NormalTok{x }\OperatorTok{\textgreater{}=}\NormalTok{ y               }\CommentTok{\# x est supérieur ou égal à y}
\NormalTok{x }\OperatorTok{\textless{}=}\NormalTok{ y               }\CommentTok{\# x est inférieur ou égal à y}
\end{Highlighting}
\end{Shaded}

\hypertarget{opuxe9rateurs-logiques}{%
\subsection{Opérateurs logiques}\label{opuxe9rateurs-logiques}}

Il existe trois opérateurs logiques : \texttt{and}, \texttt{or} et
\texttt{not}. Leur signification s\textquotesingle apparente à leur sens
en anglais. Par exemple,
\texttt{x\ \textgreater{}\ 0\ and\ x\ \textless{}\ 10} est vrai
uniquement si \texttt{x} est strictement supérieur à 0 \emph{et}
strictement inférieur à 10.

L\textquotesingle expression
\texttt{n\ \%\ 2\ ==\ 0\ or\ n\ \%\ 3\ ==\ 0} est vraie si
\emph{l\textquotesingle une} des conditions est vraie,
c\textquotesingle est-à-dire si le nombre est divisible par 2 \emph{ou}
par 3.

Enfin, l\textquotesingle opérateur \texttt{not} inverse
l\textquotesingle expression booléenne. Ainsi,
\texttt{not\ (x\ \textgreater{}\ y)} est vrai si
\texttt{x\ \textgreater{}\ y} est faux (donc si \texttt{x} est inférieur
ou égal à \texttt{y}).

\hypertarget{exuxe9cution-conditionnelle}{%
\subsection{Exécution
conditionnelle}\label{exuxe9cution-conditionnelle}}

Pour écrire des programmes utiles, nous avons presque toujours besoin de
vérifier des conditions et de modifier le comportement du programme en
conséquence. Les instructions conditionnelles (\emph{conditional
statements}) nous en donnent la capacité. La forme la plus simple est
l\textquotesingle instruction \texttt{if} :

\begin{Shaded}
\begin{Highlighting}[]
\ControlFlowTok{if}\NormalTok{ x }\OperatorTok{\textgreater{}} \DecValTok{0}\NormalTok{:}
    \BuiltInTok{print}\NormalTok{(}\StringTok{\textquotesingle{}x est positif\textquotesingle{}}\NormalTok{)}
\end{Highlighting}
\end{Shaded}

L\textquotesingle expression booléenne après le \texttt{if} est appelée
la condition. Si elle est vraie, l\textquotesingle instruction indentée
s\textquotesingle exécute. Si elle est fausse, rien ne se passe.

Les instructions \texttt{if} ont la même structure que les définitions
de fonctions : un en-tête suivi d\textquotesingle un corps indenté. Il
n\textquotesingle y a pas de limite au nombre
d\textquotesingle instructions pouvant apparaître dans le corps, mais il
doit y en avoir au moins une. Parfois, il est utile
d\textquotesingle avoir un corps sans instruction (comme espace réservé
pour du code que vous n\textquotesingle avez pas encore écrit). Dans ce
cas, vous pouvez utiliser l\textquotesingle instruction \texttt{pass},
qui ne fait rien :

\begin{Shaded}
\begin{Highlighting}[]
\ControlFlowTok{if}\NormalTok{ x }\OperatorTok{\textless{}} \DecValTok{0}\NormalTok{:}
    \ControlFlowTok{pass}  \CommentTok{\# }\AlertTok{TODO}\CommentTok{: gérer les valeurs négatives}
\end{Highlighting}
\end{Shaded}

\hypertarget{exuxe9cution-alternative}{%
\subsection{Exécution alternative}\label{exuxe9cution-alternative}}

Une deuxième forme de l\textquotesingle instruction \texttt{if} est
l\textquotesingle exécution alternative, dans laquelle il y a deux
possibilités et la condition détermine laquelle
s\textquotesingle exécute. La syntaxe ressemble à ceci :

\begin{Shaded}
\begin{Highlighting}[]
\ControlFlowTok{if}\NormalTok{ x }\OperatorTok{\%} \DecValTok{2} \OperatorTok{==} \DecValTok{0}\NormalTok{:}
    \BuiltInTok{print}\NormalTok{(}\StringTok{\textquotesingle{}x est pair\textquotesingle{}}\NormalTok{)}
\ControlFlowTok{else}\NormalTok{:}
    \BuiltInTok{print}\NormalTok{(}\StringTok{\textquotesingle{}x est impair\textquotesingle{}}\NormalTok{)}
\end{Highlighting}
\end{Shaded}

Si le reste de la division de \texttt{x} par 2 est 0, alors nous savons
que \texttt{x} est pair. Si la condition est fausse, la deuxième série
d\textquotesingle instructions s\textquotesingle exécute. Étant donné
que la condition doit être soit vraie soit fausse, exactement
l\textquotesingle une des alternatives s\textquotesingle exécutera. Ces
alternatives sont appelées des branches (\emph{branches}).

\hypertarget{conditionnelles-enchauxeenuxe9es}{%
\subsection{Conditionnelles
enchaînées}\label{conditionnelles-enchauxeenuxe9es}}

Parfois, il y a plus de deux possibilités et nous avons besoin de plus
de deux branches. Une façon d\textquotesingle exprimer un tel calcul est
d\textquotesingle utiliser une conditionnelle enchaînée (\emph{chained
conditional}) :

\begin{Shaded}
\begin{Highlighting}[]
\ControlFlowTok{if}\NormalTok{ x }\OperatorTok{\textless{}}\NormalTok{ y:}
    \BuiltInTok{print}\NormalTok{(}\StringTok{\textquotesingle{}x est inférieur à y\textquotesingle{}}\NormalTok{)}
\ControlFlowTok{elif}\NormalTok{ x }\OperatorTok{\textgreater{}}\NormalTok{ y:}
    \BuiltInTok{print}\NormalTok{(}\StringTok{\textquotesingle{}x est supérieur à y\textquotesingle{}}\NormalTok{)}
\ControlFlowTok{else}\NormalTok{:}
    \BuiltInTok{print}\NormalTok{(}\StringTok{\textquotesingle{}x et y sont égaux\textquotesingle{}}\NormalTok{)}
\end{Highlighting}
\end{Shaded}

\texttt{elif} est une abréviation de « else if ». Là encore, une seule
branche sera exécutée. Il n\textquotesingle y a pas de limite au nombre
d\textquotesingle instructions \texttt{elif}. S\textquotesingle il y a
une clause \texttt{else}, elle doit figurer à la fin, mais elle
n\textquotesingle est pas obligatoire.

Chaque condition est vérifiée dans l\textquotesingle ordre. Si la
première est fausse, la suivante est vérifiée, et ainsi de suite. Si
l\textquotesingle une d\textquotesingle elles est vraie, la branche
correspondante s\textquotesingle exécute et
l\textquotesingle instruction se termine. Même si plusieurs conditions
sont vraies, seule la première branche vraie s\textquotesingle exécute.

\hypertarget{conditionnelles-imbriquuxe9es}{%
\subsection{Conditionnelles
imbriquées}\label{conditionnelles-imbriquuxe9es}}

Une conditionnelle peut également être imbriquée dans une autre. Nous
aurions pu écrire l\textquotesingle exemple précédent de cette façon :

\begin{Shaded}
\begin{Highlighting}[]
\ControlFlowTok{if}\NormalTok{ x }\OperatorTok{==}\NormalTok{ y:}
    \BuiltInTok{print}\NormalTok{(}\StringTok{\textquotesingle{}x et y sont égaux\textquotesingle{}}\NormalTok{)}
\ControlFlowTok{else}\NormalTok{:}
    \ControlFlowTok{if}\NormalTok{ x }\OperatorTok{\textless{}}\NormalTok{ y:}
        \BuiltInTok{print}\NormalTok{(}\StringTok{\textquotesingle{}x est inférieur à y\textquotesingle{}}\NormalTok{)}
    \ControlFlowTok{else}\NormalTok{:}
        \BuiltInTok{print}\NormalTok{(}\StringTok{\textquotesingle{}x est supérieur à y\textquotesingle{}}\NormalTok{)}
\end{Highlighting}
\end{Shaded}

La conditionnelle externe contient deux branches. La première contient
une instruction simple. La deuxième branche contient une autre
instruction \texttt{if}, qui possède à son tour deux branches. Ces deux
branches sont de simples instructions, bien qu\textquotesingle elles
auraient pu être des instructions conditionnelles elles aussi.

Bien que l\textquotesingle indentation rende la structure évidente, les
conditionnelles imbriquées deviennent rapidement difficiles à lire. Il
est conseillé de les éviter quand vous le pouvez.

Les opérateurs logiques fournissent souvent un moyen de simplifier les
instructions conditionnelles imbriquées. Par exemple, nous pouvons
réécrire le code suivant en utilisant une seule condition :

\begin{Shaded}
\begin{Highlighting}[]
\ControlFlowTok{if} \DecValTok{0} \OperatorTok{\textless{}}\NormalTok{ x:}
    \ControlFlowTok{if}\NormalTok{ x }\OperatorTok{\textless{}} \DecValTok{10}\NormalTok{:}
        \BuiltInTok{print}\NormalTok{(}\StringTok{\textquotesingle{}x est un nombre positif à un chiffre.\textquotesingle{}}\NormalTok{)}
\end{Highlighting}
\end{Shaded}

L\textquotesingle instruction \texttt{print} s\textquotesingle exécute
uniquement si nous passons les deux tests, nous pouvons donc obtenir le
même effet avec l\textquotesingle opérateur \texttt{and} :

\begin{Shaded}
\begin{Highlighting}[]
\ControlFlowTok{if} \DecValTok{0} \OperatorTok{\textless{}}\NormalTok{ x }\KeywordTok{and}\NormalTok{ x }\OperatorTok{\textless{}} \DecValTok{10}\NormalTok{:}
    \BuiltInTok{print}\NormalTok{(}\StringTok{\textquotesingle{}x est un nombre positif à un chiffre.\textquotesingle{}}\NormalTok{)}
\end{Highlighting}
\end{Shaded}

\hypertarget{ruxe9cursivituxe9}{%
\subsection{Récursivité}\label{ruxe9cursivituxe9}}

Il est légal pour une fonction d\textquotesingle appeler une autre
fonction ; il est également légal pour une fonction de
s\textquotesingle appeler elle-même. L\textquotesingle exemple suivant
illustre ce concept, appelé récursivité (\emph{recursion}) :

\begin{Shaded}
\begin{Highlighting}[]
\KeywordTok{def}\NormalTok{ compte\_a\_rebours(n):}
    \ControlFlowTok{if}\NormalTok{ n }\OperatorTok{\textless{}=} \DecValTok{0}\NormalTok{:}
        \BuiltInTok{print}\NormalTok{(}\StringTok{\textquotesingle{}Décollage !\textquotesingle{}}\NormalTok{)}
    \ControlFlowTok{else}\NormalTok{:}
        \BuiltInTok{print}\NormalTok{(n)}
\NormalTok{        compte\_a\_rebours(n }\OperatorTok{{-}} \DecValTok{1}\NormalTok{)}
\end{Highlighting}
\end{Shaded}

Si \texttt{n} est inférieur ou égal à 0, la fonction affiche le mot «
Décollage ! ». Sinon, elle affiche \texttt{n}, puis appelle une fonction
nommée \texttt{compte\_a\_rebours} (elle-même), en lui passant
\texttt{n\ -\ 1} comme argument.

Une fonction qui s\textquotesingle appelle elle-même est dite récursive
; le processus d\textquotesingle exécution de cette fonction
s\textquotesingle appelle la récursivité.

\hypertarget{diagrammes-de-pile-pour-les-fonctions-ruxe9cursives}{%
\subsection{Diagrammes de pile pour les fonctions
récursives}\label{diagrammes-de-pile-pour-les-fonctions-ruxe9cursives}}

Au chapitre précédent, nous avons utilisé un diagramme de pile pour
représenter l\textquotesingle état d\textquotesingle un programme
pendant un appel de fonction. Le même type de diagramme peut aider à
interpréter une fonction récursive.

À chaque fois qu\textquotesingle une fonction est appelée, Python crée
un nouveau contexte local (un cadre ou \emph{frame}), qui contient les
variables et paramètres locaux de la fonction. Pour une fonction
récursive, il peut y avoir plusieurs cadres simultanément sur la pile,
chacun correspondant à un appel différent de la fonction, avec sa propre
valeur de \texttt{n}.

\hypertarget{ruxe9cursivituxe9-infinie}{%
\subsection{Récursivité infinie}\label{ruxe9cursivituxe9-infinie}}

Si une récursivité n\textquotesingle atteint jamais de cas de base, elle
s\textquotesingle exécute pour toujours, et le programme ne se termine
jamais. C\textquotesingle est ce qu\textquotesingle on appelle une
récursivité infinie (\emph{infinite recursion}).

Dans la plupart des environnements de programmation, un programme avec
une récursivité infinie ne s\textquotesingle exécute pas réellement pour
toujours. Python signale un message d\textquotesingle erreur lorsque la
profondeur de récursivité maximale est atteinte (généralement une
\texttt{RecursionError}).

\hypertarget{saisie-au-clavier}{%
\subsection{Saisie au clavier}\label{saisie-au-clavier}}

Les programmes que nous avons écrits jusqu\textquotesingle à présent
n\textquotesingle acceptent aucune saisie de
l\textquotesingle utilisateur. Ils effectuent le même calcul à chaque
fois. Python fournit une fonction intégrée appelée \texttt{input} qui
met l\textquotesingle exécution du programme en pause et attend que
l\textquotesingle utilisateur tape quelque chose.

\begin{Shaded}
\begin{Highlighting}[]
\OperatorTok{\textgreater{}\textgreater{}\textgreater{}}\NormalTok{ texte }\OperatorTok{=} \BuiltInTok{input}\NormalTok{()}
\NormalTok{Quelque chose}
\OperatorTok{\textgreater{}\textgreater{}\textgreater{}} \BuiltInTok{print}\NormalTok{(texte)}
\NormalTok{Quelque chose}
\end{Highlighting}
\end{Shaded}

Avant d\textquotesingle obtenir une saisie, il est judicieux
d\textquotesingle afficher une invite indiquant à
l\textquotesingle utilisateur ce qu\textquotesingle il doit entrer. Vous
pouvez passer une chaîne de caractères à \texttt{input} pour
qu\textquotesingle elle s\textquotesingle affiche avant la mise en pause
:

\begin{Shaded}
\begin{Highlighting}[]
\OperatorTok{\textgreater{}\textgreater{}\textgreater{}}\NormalTok{ nom }\OperatorTok{=} \BuiltInTok{input}\NormalTok{(}\StringTok{\textquotesingle{}Quel est votre nom ?}\CharTok{\textbackslash{}n}\StringTok{\textquotesingle{}}\NormalTok{)}
\NormalTok{Quel est votre nom ?}
\NormalTok{Arthur}
\OperatorTok{\textgreater{}\textgreater{}\textgreater{}} \BuiltInTok{print}\NormalTok{(nom)}
\NormalTok{Arthur}
\end{Highlighting}
\end{Shaded}

La séquence \texttt{\textbackslash{}n} à la fin de la chaîne représente
une nouvelle ligne, ce qui place l\textquotesingle entrée de
l\textquotesingle utilisateur sur la ligne en dessous de
l\textquotesingle invite.

Si vous attendez un entier, vous pouvez essayer de convertir la valeur
de retour avec \texttt{int()} :

\begin{Shaded}
\begin{Highlighting}[]
\OperatorTok{\textgreater{}\textgreater{}\textgreater{}}\NormalTok{ invite }\OperatorTok{=} \StringTok{\textquotesingle{}Quelle est la vitesse de vol d}\CharTok{\textbackslash{}\textquotesingle{}}\StringTok{une hirondelle à vide ?}\CharTok{\textbackslash{}n}\StringTok{\textquotesingle{}}
\OperatorTok{\textgreater{}\textgreater{}\textgreater{}}\NormalTok{ vitesse }\OperatorTok{=} \BuiltInTok{input}\NormalTok{(invite)}
\NormalTok{Quelle est la vitesse de vol d}\StringTok{\textquotesingle{}une hirondelle à vide ?}
\ErrorTok{42}
\ErrorTok{\textgreater{}\textgreater{}\textgreater{} int}\NormalTok{(vitesse)}
\DecValTok{42}
\end{Highlighting}
\end{Shaded}

\hypertarget{duxe9bogage}{%
\subsection{Débogage}\label{duxe9bogage}}

Les messages d\textquotesingle erreur et les traces
d\textquotesingle exécution (\emph{tracebacks}) de Python indiquent où
s\textquotesingle est produite l\textquotesingle erreur, mais ne disent
généralement pas pourquoi. Si une fonction s\textquotesingle appelle
elle-même trop souvent, la trace peut devenir extrêmement longue et
difficile à lire.

Lorsque vous écrivez des conditions complexes, assurez-vous de bien
tester les différents chemins d\textquotesingle exécution (les cas
extrêmes, les nombres négatifs, les zéros). Ajouter des instructions
\texttt{print} temporaires peut grandement faciliter la compréhension du
chemin réellement emprunté par votre programme.

\hypertarget{glossaire}{%
\subsection{Glossaire}\label{glossaire}}

opérateur modulo modulus operator Un opérateur, noté avec un signe de
pourcentage (\texttt{\%}), qui opère sur des entiers et donne le reste
lorsqu\textquotesingle un nombre est divisé par un autre.

expression booléenne boolean expression Une expression dont la valeur
est soit \texttt{True}, soit \texttt{False}.

opérateur relationnel relational operator L\textquotesingle un des
opérateurs qui compare ses opérandes : \texttt{==}, \texttt{!=},
\texttt{\textgreater{}}, \texttt{\textless{}}, \texttt{\textgreater{}=}
et \texttt{\textless{}=}.

opérateur logique logical operator L\textquotesingle un des opérateurs
qui combine des expressions booléennes : \texttt{and}, \texttt{or} et
\texttt{not}.

instruction conditionnelle conditional statement Une instruction qui
contrôle le flux d\textquotesingle exécution en fonction
d\textquotesingle une condition.

\begin{description}
\item[condition]
L\textquotesingle expression booléenne dans une instruction
conditionnelle qui détermine quelle branche s\textquotesingle exécute.
\end{description}

branche branch L\textquotesingle une des séquences alternatives
d\textquotesingle instructions dans une conditionnelle.

conditionnelle enchaînée chained conditional Une instruction
conditionnelle comportant une série d\textquotesingle alternatives
(\texttt{elif}).

conditionnelle imbriquée nested conditional Une instruction
conditionnelle qui apparaît dans l\textquotesingle une des branches
d\textquotesingle une autre instruction conditionnelle.

récursivité recursion Le processus d\textquotesingle appel de la
fonction qui est elle-même en cours d\textquotesingle exécution.

cas de base base case Une branche conditionnelle dans une fonction
récursive qui n\textquotesingle effectue pas d\textquotesingle appel
récursif.

récursivité infinie infinite recursion Une récursivité qui
n\textquotesingle a pas de cas de base, ou ne l\textquotesingle atteint
jamais. Finit par provoquer une erreur d\textquotesingle exécution.

\hypertarget{exercices}{%
\subsection{Exercices}\label{exercices}}

\textbf{Exercice 1}

Le dernier théorème de Fermat stipule qu\textquotesingle il
n\textquotesingle existe aucun entier strictement positif \emph{a},
\emph{b} et \emph{c} tel que :

\emph{a}\textsuperscript{n} + \emph{b}\textsuperscript{n} =
\emph{c}\textsuperscript{n}

pour toutes valeurs de \emph{n} strictement supérieures à 2.

\begin{enumerate}
\def\labelenumi{\arabic{enumi}.}
\tightlist
\item
  Écrivez une fonction nommée \texttt{verifier\_fermat} qui prend quatre
  paramètres \texttt{a}, \texttt{b}, \texttt{c} et \texttt{n}, et qui
  vérifie si le théorème de Fermat est vrai. Si \emph{n} est plus grand
  que 2 et que l\textquotesingle équation est vraie, le programme doit
  afficher « Bon sang, Fermat avait tort ! ». Sinon, le programme doit
  afficher « Non, cela ne fonctionne pas. »
\item
  Écrivez une fonction qui invite l\textquotesingle utilisateur à saisir
  des valeurs pour \texttt{a}, \texttt{b}, \texttt{c} et \texttt{n}, les
  convertit en entiers, et utilise \texttt{verifier\_fermat} pour
  vérifier si elles violent le théorème de Fermat.
\end{enumerate}

\textbf{Exercice 2}

Si l\textquotesingle on vous donne trois bâtons, vous pouvez ou non être
en mesure de les disposer en triangle. Par exemple, si
l\textquotesingle un des bâtons a une longueur de 12 pouces et que les
deux autres font 1 pouce chacun, vous ne pourrez pas faire se rencontrer
les petits bâtons au milieu. Pour tout ensemble de trois longueurs, il
existe un test simple pour vérifier s\textquotesingle il est possible de
former un triangle :

\emph{Si l\textquotesingle une des trois longueurs est supérieure à la
somme des deux autres, alors vous ne pouvez pas former de triangle.
Sinon, vous le pouvez.}

\begin{enumerate}
\def\labelenumi{\arabic{enumi}.}
\tightlist
\item
  Écrivez une fonction nommée \texttt{est\_un\_triangle} qui prend trois
  entiers comme arguments, et qui affiche « Oui » ou « Non », selon que
  vous pouvez ou non former un triangle à partir de bâtons possédant les
  longueurs données.
\item
  Écrivez une fonction qui invite l\textquotesingle utilisateur à saisir
  trois longueurs de bâton, les convertit en entiers, et utilise
  \texttt{est\_un\_triangle} pour vérifier si des bâtons de ces
  longueurs peuvent former un triangle.
\end{enumerate}
